\chapter{Interacción bidireccional: conjugación de propagadores avanzado y retardado y condiciones de borde del absorbedor}
\label{chap:mt-accion-interaccion-absorbedor}

Una ruta se vuelve dinámica cuando recibe un funcional de acción y unas
condiciones de frontera. La acción asigna peso a la historia completa; su
variación selecciona transportes compatibles con los extremos; la
interacción acopla grados de libertad que comparten una frontera; la lectura
bidireccional conserva las orientaciones temporal directa e inversa del
mismo antecedente. En este marco, emisión y absorción forman extremos de un
problema de contorno común.

El absorbedor ocupa una posición precisa en la arquitectura. HMT aporta la
banda orientada, la memoria de ida y retorno y las involuciones discretas.
La Mecánica Dimensional añade el operador de evolución, las corrientes y el
transductor de acción. Una realización física de absorbedor enlaza esa
estructura con propagadores avanzados y retardados y con condiciones de
borde causales. Las identidades algebraicas del primer tramo están
demostradas en sus dominios; el último enlace mantiene el estatuto de
interfaz física hasta construir el entrelazador completo.

\section{Funcional de acción sobre rutas enriquecidas}

Sea $\gamma=((c_0,\mu_0),(d_1,\ldots,d_L))$ una ruta finita, con direcciones
$d_t$ y hojas $\mu_t\in\{\Sigma,\Pi\}$. Definimos
\[
 \operatorname{turn}(\gamma)
 =\#\{2\leq t\leq L:d_t\neq d_{t-1}\},
\]
\[
 \operatorname{switch}(\gamma)
 =\#\{2\leq t\leq L:\mu_t\neq\mu_{t-1}\}.
\]

\begin{definicionnarrativa}[title={Acción discreta de interacción mínima}]
\label{def:mt-accion-minima}
Para $\lambda\geq0$, la acción local de ruta es
\begin{equation}
 \mathcal A_\lambda(\gamma)
 =\operatorname{turn}(\gamma)
 +\lambda\operatorname{switch}(\gamma)
 =\sum_{t=2}^{L}
 \bigl[1-\delta_{d_t,d_{t-1}}
 +\lambda(1-\delta_{\mu_t,\mu_{t-1}})\bigr].
 \label{eq:mt-accion-minima}
\end{equation}
La primera contribución mide variación direccional y la segunda mide
transporte entre hojas.
\end{definicionnarrativa}

La acción general incorpora también borde, acarreo y memoria:
\begin{equation}
 \mathcal S[\gamma]
 =\sum_{e\in\gamma}\mathcal L_e
 +\mathcal B_{\rm in}(\partial_-\gamma)
 +\mathcal B_{\rm out}(\partial_+\gamma)
 +\mathcal M(\gamma).
 \label{eq:mt-accion-ruta-general}
\end{equation}
Aquí $\mathcal L_e$ es el coste local, $\mathcal B_{\rm in/out}$ fija los
datos admisibles de los extremos y $\mathcal M$ conserva la información que
la publicación visible contrae. La variación se realiza dentro de una clase
de rutas con los mismos tipos de frontera.

\begin{lecturafisica}[title={Principio de mínima interacción compatible con memoria}]
\label{prin:mt-minima-interaccion}
Entre las prolongaciones compatibles con los datos de frontera, la
realización física seleccionada extrema el funcional de acción dentro del
sector que conserva la memoria exigida por la composición.
\end{lecturafisica}

El principio fija la clase del problema variacional. Una dinámica concreta
añade el Hamiltoniano o lagrangiano del sector y demuestra que su extremo
produce las ecuaciones de movimiento. La memoria funciona como restricción
dinámica: dos rutas con la misma sombra final pueden pertenecer a clases
variacionales distintas.

\section{Interacción como acoplamiento tipado}

Sean $\mathfrak p_1$ y $\mathfrak p_2$ dos secciones realizadas sobre rutas
$\gamma_1$ y $\gamma_2$. Una interacción en una frontera $B$ consiste en
un morfismo
\begin{equation}
 \mathcal I_B:
 E^{(1)}_B\otimes E^{(2)}_B
 \longrightarrow E^{(12)}_B
 \label{eq:mt-interaccion-frontera}
\end{equation}
y en un término de acción $\mathcal S_{\rm int}$ compatible con las
representaciones del sector. El morfismo conserva los tipos compartidos y
transforma los grados internos según las reglas de acoplamiento. La
composición completa adopta la forma
\[
 (\gamma_2,s_2)\star_{(B,\mathcal I_B)}(\gamma_1,s_1),
\]
donde el registro de $B$ se cuenta una vez y la memoria de los dos lados se
transporta al estado compuesto.

La covariancia de la interacción exige, para cada transformación de simetría
$g$,
\begin{equation}
 \rho_{12}(g)\,\mathcal I_B
 =\mathcal I_B\bigl(\rho_1(g)\otimes\rho_2(g)\bigr).
 \label{eq:mt-covariancia-interaccion}
\end{equation}
Esta identidad decide si el acoplamiento pertenece a la representación
física declarada. En los sectores electromagnético y de color, la misma
estructura se expresa mediante una conexión y su transporte paralelo.

\section{Conjugación temporal del propagador}

Sea $\mathcal H_{\mathbb R}$ la realificación del espacio cuántico, dotada
de una estructura compleja $J_E$ con $J_E^2=-I$. Sea $H=H^*$ un operador
autoadjunto con dominio común invariante y sea $C$ una involución ortogonal
que satisface
\begin{equation}
 HJ_E=J_EH,\qquad
 CJ_EC^{-1}=-J_E,
 \qquad CHC^{-1}=H.
 \label{eq:mt-hipotesis-reversion-propagador}
\end{equation}

\begin{resultado}[title={Teorema: Reversión bidireccional del grupo de evolución}]
\label{thm:mt-reversion-propagador}
El operador
\begin{equation}
 U(t)=\exp\!\left(-\frac{J_EHt}{\hbar_{\rm int}}\right)
 \label{eq:mt-grupo-evolucion}
\end{equation}
forma un grupo unitario uniparamétrico y cumple
\begin{equation}
 \boxed{CU(t)C^{-1}=U(-t)}.
 \label{eq:mt-reversion-unitaria}
\end{equation}
\end{resultado}

\paragraph{Demostración de la reversión bidireccional del grupo unitario \(U(t)\): \(CU(t)C^{-1}=U(-t)\)}
El generador $A=-J_EH/\hbar_{\rm int}$ es antiadjunto. El teorema de Stone
produce el grupo unitario $U(t)=e^{tA}$. Las hipótesis de
\eqref{eq:mt-hipotesis-reversion-propagador} dan $CAC^{-1}=-A$; el cálculo
funcional entrega $Ce^{tA}C^{-1}=e^{-tA}$.


La involución $C_M=-\operatorname{rev}$ de las palabras de ruta aporta el
dato genealógico de orientación. Cuando el funtor físico lleva $C_M$ a la
involución $C$ del teorema, las secciones partícula y antipartícula quedan
representadas por orientaciones temporales conjugadas del mismo grupo
unitario. La identidad está demostrada con la estructura de partida
explícita: autoadjunción, dominio, estructura compleja e involución.

\section{Propagadores avanzados y retardados}

Sea $L$ el operador hiperbólico de una realización de campos y sean
$G_{\rm ret}$ y $G_{\rm adv}$ sus inversos fundamentales con soportes
causales futuro y pasado, respectivamente:
\begin{equation}
 LG_{\rm ret}=LG_{\rm adv}=I,
 \qquad
 \operatorname{supp}(G_{\rm ret}f)\subset J^+(\operatorname{supp}f),
 \qquad
 \operatorname{supp}(G_{\rm adv}f)\subset J^-(\operatorname{supp}f).
 \label{eq:mt-green-avanzado-retardado}
\end{equation}
La diferencia causal
$E=G_{\rm ret}-G_{\rm adv}$ transporta datos entre soluciones y la media
simétrica
\[
 G_{\rm sym}=\frac12(G_{\rm ret}+G_{\rm adv})
\]
realiza un núcleo temporalmente simétrico.

Una interfaz HMT--MD hacia esta teoría es una familia de mapas
\begin{equation}
 \mathfrak I_{\rm prop}:
 (\gamma,\gamma^\vee,\mathcal B_\gamma,U,C_M)
 \longrightarrow
 (G_{\rm ret},G_{\rm adv},E,\mathcal H_{\rm field})
 \label{eq:mt-interfaz-green}
\end{equation}
que entrelaza la involución de ruta con el intercambio
$G_{\rm ret}\leftrightarrow G_{\rm adv}$, preserva el producto interno o la
forma simpléctica pertinente y lleva las acciones de ambos lados a la misma
forma funcional. Este mapa constituye la condición matemática que convierte
la orientación bidireccional interna en propagación avanzada--retardada de
campos.

\section{Realización HMT–MD del absorbedor mediante inversión de rutas y conjugación partícula–antipartícula}

\begin{definicionnarrativa}[title={Sistema de emisión--absorción bidireccional}]
\label{def:mt-absorbedor}
Una realización física de absorbedor es un problema de contorno formado por
una corriente emisora $J_{\rm em}$, una respuesta absorbente $J_{\rm abs}$,
los propagadores $G_{\rm ret}$ y $G_{\rm adv}$, y condiciones de borde que
cierran conjuntamente las dos orientaciones del transporte. Su acción
efectiva temporalmente simétrica posee la forma
\begin{equation}
 \mathcal S_{\rm abs}[J]
 =\mathcal S_{\rm matter}[J]
 +\frac12\langle J,G_{\rm sym}J\rangle,
 \qquad J=J_{\rm em}+J_{\rm abs},
 \label{eq:mt-accion-absorbedor}
\end{equation}
y sus condiciones de contorno determinan la parte radiativa observada.
\end{definicionnarrativa}

La definición expresa una teoría de acción a distancia o de respuesta de
contorno en lenguaje operatorio. Dentro de HMT--MD, la banda aporta el par
$(\gamma,\gamma^\vee)$; el resultado de reversión bidireccional aporta
la reversión exacta del grupo unitario; el funtor
$\mathfrak I_{\rm prop}$ debe aportar causalidad, producto interno y
condiciones de borde. La arquitectura del absorbedor prolonga la estructura
bidireccional ya construida. Su realización permanece como interfaz física,
y la construcción del entrelazador causal completo constituye un problema
matemático explícitamente delimitado.

El cierre físico exige cuatro identidades verificables:
\begin{enumerate}
 \item conmutatividad entre conjugación de ruta e intercambio de los dos
 propagadores;
 \item igualdad de las acciones transportadas por
 $\mathfrak I_{\rm prop}$;
 \item conservación de la forma simpléctica o del producto interno;
 \item compatibilidad de las condiciones de emisión y absorción bajo
 refinamiento de la ruta.
\end{enumerate}
Un fallo en cualquiera de estas identidades refuta la realización particular
y deja intacta la reversión operatoria ya demostrada en
\eqref{eq:mt-reversion-unitaria}.

\section{Las acciones separadas de \texorpdfstring{$C$}{C},
\texorpdfstring{$P$}{P} y \texorpdfstring{$T$}{T}}

Sobre una ruta completa definimos tres involuciones tipadas:
\begin{enumerate}
 \item $C$ intercambia las hojas $\Sigma\leftrightarrow\Pi$ y conjuga los
 caracteres orientados;
 \item $P$ refleja la carta espacial, por ejemplo
 $(i,j)\mapsto(i,-j)\pmod9$, junto con
 $E\leftrightarrow W$;
 \item $T$ invierte el orden del registro, reemplaza cada dirección por su
 opuesta y permuta los extremos de la ruta.
\end{enumerate}
Cada operación actúa sobre un componente distinto del estado. Su composición
$CPT$ actúa sobre hoja, espacio y orientación temporal a la vez.

\begin{resultado}[title={Teorema: Invariancia discreta individual y conjunta}]
\label{thm:mt-cpt-rutas}
Para toda ruta $\gamma$ y todo $\lambda\geq0$,
\begin{equation}
 \mathcal A_\lambda(\gamma)
 =\mathcal A_\lambda(C\gamma)
 =\mathcal A_\lambda(P\gamma)
 =\mathcal A_\lambda(T\gamma)
 =\mathcal A_\lambda(CPT\gamma).
 \label{eq:mt-invariancia-cpt-rutas}
\end{equation}
\end{resultado}

\paragraph{Compatibilidad de las acciones \(C\), \(P\) y \(T\) con la reversión unitaria \(U(t)\mapsto U(-t)\)}
$P$ permuta el alfabeto de direcciones y conserva el patrón de igualdades
entre pasos consecutivos. $C$ permuta globalmente las hojas y conserva el
patrón de cambios de hoja. $T$ invierte el orden de ambas sucesiones y
conserva el número de transiciones adyacentes. Cada involución preserva por
separado los dos sumandos de
\eqref{eq:mt-accion-minima}; la composición hereda la invariancia.


El teorema pertenece al registro discreto y su acción local. La matriz CKM
realizada en la base de sabor posee un invariante de Jarlskog distinto de
cero y, por tanto, violación CP en esa representación. Ambos resultados son
compatibles: un lector orientado puede distinguir estados relacionados por
$CP$ mientras el funcional local
$\mathcal A_\lambda$ conserva cada involución. La
realización del teorema CPT relativista de campos añade localidad,
covariancia de Lorentz, estructura espectral y representación de campos. En
ese dominio ampliado, CPT conjunta constituye la simetría a preservar. En el
dominio interno, el cierre CKM y la expansión temporal permanecen compatibles
con la invariancia CPT demostrada para la acción discreta.

Quedan así separados tres estratos. La invariancia de
\eqref{eq:mt-invariancia-cpt-rutas} es un teorema interno exacto. La identidad
\eqref{eq:mt-reversion-unitaria} está demostrada con estructura de partida
explícita. La teoría física del absorbedor y el teorema CPT de campos exigen
los entrelazadores propios de sus codominios. Esta separación mantiene la
orientación temporal inversa dentro de la definición de partícula y concede
a cada afirmación la fuerza que le corresponde.
